\section{Source audit, scope corrections, and integration notes}
\label{sec:audit}

An audit is useful only when it distinguishes an actual error from a theorem
that simply does not answer a new question. The inspected bridge chapters
already contain many of the corrections described in earlier incoming
reports. We found no additional substantive mathematical error in the
specific baseline theorems used as inputs here. The following observations
identify the boundaries that must be preserved when integrating this work.

\subsection{A verified normalization typo in an external source}
There is one concrete bibliographical correction. In the published PDF of
Gay--Sebbar \cite{gamma:GaySebbar2024}, printed page~4, the unnumbered
display immediately after Eq.~(1) and before Eq.~(2) omits the factor $z$
between the Lerch transcendent and the polylogarithm. With that paper's
preceding zero-based definition, the corrected display is
\begin{equation}\label{audit:lerch-typo}
 \Li_s(z)=z\Phi(z,s,1)
      =\sum_{n=1}^{\infty}\frac{z^n}{n^s}.
\end{equation}
Indeed multiplication by $z$ and the reindexing $m=n+1$ prove the identity
immediately in $|z|<1$. The omission was checked in the published PDF, so
it is not merely an HTML transcription issue. The inspected ProveIt source
already has the correct convention: \path{chapters/09-zero-geometry.tex},
label \texttt{zeros:eq:lambda} and its subsequent $a=1$ specialization,
includes the corresponding factor $1/\rho$. This is an external-source
typographical correction, not a required change to that manuscript formula.

\subsection{What the present results settle}
The final research paragraph of
\path{chapters/09-herglotz-jets.tex} asks for a transition theory in which
the argument varies with derivative order \cite{repoHerglotz}.
Theorem~\ref{herg:thm:expansion} answers that question on every compact
positive interval of $x/r$, with all fixed ratio derivatives and a complete
inverse-order expansion. Theorem~\ref{herg:thm:mixture} also supplies a
finite-order inequality on the full positive axis.
Corollary~\ref{herg:cor:ratios} identifies the leading limit along every
extended limiting ratio. A uniformly accurate expansion that also resolves
the exponentially small correction as the ratio approaches zero remains
a further problem; the compact-ratio theorem does not assert it.

The multivariate discussion in
\path{chapters/08-integral-jets.tex} correctly retains the full Koszul
complex outside the one-variable case \cite{repoJets}. Theorem~
\ref{ns:thm:two-generator} gives a new sufficient condition for a closed
dimension law: local Frobenius duality and an active ideal needing at most
two generators. In particular, it settles the dimension and abelian Smith
factors for every active sequence with two entries over a rectangular jet
algebra. It does not classify all multivariate torsion modules.

The three-generator example in Proposition~\ref{ns:prop:counterexample}
is a counterexample to an \emph{unrestricted extrapolation}, not to the
manuscript's properly restricted theorem. Similarly,
\eqref{ns:eq:module-obstruction} disproves a stronger module substitution
that is not needed for the valid dimension formula. These distinctions
should remain explicit in the canonical manuscript.

\subsection{Normalization and domain statements to preserve}
\begin{enumerate}
\item \textbf{Lerch normalization.}
With $\Phi(z,s,a)=\sum_{n\geq0}z^n(n+a)^{-s}$, one has
$\Li_s(z)=z\Phi(z,s,1)$. Accordingly the exponentially shifted function
in \eqref{res:F} equals $\Li_s(e^{-t})$ at $a=1$.
Dropping the factor $e^{-at}$ would change the exact finite-part expansion.
The Laurent sign convention \eqref{scope:laurent} must accompany every
formula for generalized Stieltjes constants.

\item \textbf{Uniformity and limits.}
The coefficients $b_{k,j}$ in Theorem~\ref{res:main} retain their dependence
on $k$. Their limiting values from Proposition~\ref{res:highorder} cannot
replace them in an arbitrary all-order joint expansion without a rate
condition linking $k$ and $T$. Similarly, a fixed-argument Herglotz
asymptotic cannot be substituted at $x=r\lambda$ without a new remainder
estimate. The exact mixture provides that estimate here.

\item \textbf{Antiderivative constants.}
The balanced functions in \eqref{gamma:eq:balanced} have zero mean and,
from the second level onward, matching periodic endpoints. A primitive
normalized to vanish at a chosen base point usually differs by a polynomial
after repeated integration. Correlation identities for the balanced family
cannot be transferred unchanged to that other normalization.

\item \textbf{Ordinary integrals versus regularized distributions.}
At a nonzero shift, the two Hurwitz endpoint singularities are separate and
the domain is $\Re s,\Re t<1$. At zero shift the additional condition
$\Re(s+t)<1$ is essential. Argument differentiation to the digamma level
produces a $1/x$ singularity, so an unqualified ordinary-integral extension
would be invalid. A finite-part or distributional convention must be stated
and its possible endpoint terms proved.

\item \textbf{Which function has the imaginary zeros.}
Theorem~\ref{herg:thm:profilezeros} classifies the zeros of
$H(y)-\zeta(2)/2$. It does not classify the zeros of $H$ itself. The
centering is what produces the odd meromorphic function and the positive
Stieltjes-transform argument.

\item \textbf{Formal rank versus values.}
The rank and torsion results concern the distribution presentation
\eqref{ns:eq:distribution} over its stated coefficient ring. They do not
prove linear or algebraic independence of evaluated polylogarithms,
Hurwitz jets, generalized Stieltjes constants, or Gamma values. If an
arithmetic specialization inverts two, the reflected two-torsion disappears;
the choice of lattice therefore matters.
\end{enumerate}

\subsection{Numerical pitfalls and the accompanying checks}
Two implementation issues appeared during independent verification and were
corrected in the accompanying code. These concern evaluation methods, not
counterexamples to the analytic identities.

First, direct differentiation in the order of a generic polylogarithm
implementation near an integer can evaluate two large pole terms before
their cancellation. An infinitesimal finite-difference step may then lose
many digits even at apparently generous working precision. The correlation
checker uses a resolved fourth-order symmetric stencil with additional
precision, and cross-checks it against independent Hurwitz-jet formulas and
split endpoint quadrature. The reusable correlation module uses the
convergent odd-zeta expansion instead. The resonance checker pairs the
poles analytically and uses native Hurwitz order derivatives for its
finite-part comparison.

Second, the high-derivative Herglotz formula multiplies small Hurwitz
differences by large powers. Its checker adds guard precision depending
on both order and argument before this multiplication. Working precision
alone is not an error certificate. Every analytic tail estimate in this
article bounds omitted terms in exact arithmetic; a complete implementation
with rigorous rounding bounds would additionally require interval or ball
arithmetic.

\begin{center}
\begin{tabular}{@{}>{\raggedright\arraybackslash}p{0.31\textwidth}>{\raggedright\arraybackslash}p{0.59\textwidth}@{}}
\toprule
Check & Independent evidence recorded in the package\\
\midrule
Resonant cancellation & Six direct Lerch comparisons at 100 decimal
working digits; largest absolute residual below $3.4\times10^{-85}$\\[4pt]
Stieltjes finite parts & Orders $0,1,2$ checked against separately summed
Hurwitz derivative tails; largest residual below $1.6\times10^{-67}$\\[4pt]
Shifted Hurwitz products & Split quadrature against the polylogarithmic
kernel for real and complex orders; residuals below $6.1\times10^{-33}$\\[4pt]
Log-Gamma correlation & Independent quadrature, Hurwitz jets, polylogarithm
jets, and the convergent odd-zeta expansion at three interior shifts\\[4pt]
Herglotz transition & Independent Hurwitz evaluations, exact central-moment
coefficients, finite Taylor inequalities, and Bessel/primitive comparisons\\[4pt]
Formal reflection jets & Exact binary elimination in the original point
presentation and in separate Koszul matrices; dimensions $44,210,232$\\
\bottomrule
\end{tabular}
\end{center}

The numerical rows are diagnostics. The exact rank rows are finite proofs
of their particular matrix statements. The general theorems rely on the
proofs in the body of the article. The data files label asymptotic
differences separately from identity residuals; a truncated asymptotic
formula is not expected to agree to working precision.

\subsection{Suggested canonical placement}
The source files are modular so that the continuation can be integrated
without importing a second preamble. The following placement respects the
manuscript's existing organization.
\begin{center}
\begin{tabular}{@{}>{\raggedright\arraybackslash}p{0.44\textwidth}>{\raggedright\arraybackslash}p{0.47\textwidth}@{}}
\toprule
Contribution & Suggested destination\\
\midrule
Uniform integer-order resonance and finite parts & Lerch boundary
continuation; cross-reference the zeta-jet and Stieltjes sections\\[4pt]
Shifted products, convexity, and balanced primitives & Chapter 07,
adjacent to the existing integration and log-Gamma calculus\\[4pt]
Herglotz transition, profile, and antiderivatives & Chapter 09,
after the fixed-argument derivative asymptotic\\[4pt]
Two-generator theorem and obstruction & Chapter 08,
after the finite-free Koszul and multivariate jet discussion\\
\bottomrule
\end{tabular}
\end{center}

The accompanying \path{INTEGRATION.md} gives file-level instructions and
the theorem labels used here. The source manifest records the frozen
repository revision and the locally reviewed source files. The archive
contains the complete article, reproducible code, recorded outputs, and
figure source. The research contribution is mathematical prose with
independent exact and numerical checks; it is not a Lean formalization.

